% !TeX program = xelatex
% Source authority: Noether_German_Cumulative_v26_R821integrated_20260716.tex
% Covered source segments: P06-S0002, P06-S0004, P06-S0005
% Status: working Arabic translation; external/native mathematical review pending
\documentclass[12pt]{article}
\usepackage[a4paper,margin=27mm]{geometry}
\usepackage{fontspec}
\usepackage{polyglossia}
\setdefaultlanguage{arabic}
\setotherlanguage{english}
\newfontfamily\arabicfont[Script=Arabic,Scale=1.08]{Arial}
\newfontfamily\englishfont{TeX Gyre Termes}
\setlength{\parindent}{1.5em}
\setlength{\parskip}{0.65em}
\emergencystretch=2em

\begin{document}

\begin{center}
  {\Large\bfseries ٦. الحقول وأنظمة الدوال النسبية\par}
  \vspace{0.8em}
  \textenglish{Math. Ann. 76 (1915), pp. 161--196}
\end{center}

\vspace{0.5em}
{\footnotesize\noindent
ترجمة عمل عربية، الإصدار ٠٠١. المصدر الألماني المعتمد مؤرخ في
\textenglish{2026-07-16}. المقاطع:
\textenglish{P06-S0002, P06-S0004, P06-S0005}.
تنتظر هذه الصياغة مراجعةً رياضيةً ولغويةً خارجية.\par}
\vspace{0.5em}
\hrule
\vspace{1.2em}

يتناول هذا البحث مسائل الأساس في الأنظمة الاعتباطية من الدوال النسبية
والدوال النسبية الصحيحة؛ وتُبيّن طرائق نظرية الحقول المستخدمة أن لبّ هذه
المسائل هو حلّها للحقول المؤلفة من دوال نسبية --- أي حقول الدوال
النسبية\footnote{للمقارنة، انظر البلاغ التمهيدي الذي قُدّم بمناسبة اجتماع
علماء الطبيعة في فيينا: «حقول الدوال النسبية»،
\textenglish{Jahresber. d. D. Math.-Ver. 22 (1913)}، حيث عُرضت نظرة عامة
إلى المسائل المطروحة والنتائج المتعلقة بحقول الدوال.} ---، في حين ينتج
تعميم النتائج على الأنظمة الاعتباطية بوصفه نتيجةً لذلك.

ومن مسائل الأساس في الأنظمة العامة، لم يُعرف حتى الآن إلا وجود أساس
المودول الذي تضمنه مبرهنة هيلبرت (\textenglish{Math. Ann. 36}). وفيما يلي
تُجاب مسألة قابلية التمثيل النسبي إجابةً كاملة بإثبات وجود أساس نسبي لكل
نظام اعتباطي (\textenglish{\S~7})؛ أما الأساس النسبي للحقول فيُستخرج منذ
(\textenglish{\S~4}). ويتيح وجود هذا الأساس النسبي أن ننطلق في جميع
المواضع من الحقل أو النظام المعرّف تعريفًا مجردًا، وأن نتجنب بذلك الصعوبات
التي لا يسببها النظام ذاته بل الاختيار الخاص للأساس النسبي فحسب، مثل ظهور
مقامات خاصة أو نقاط أساسية لدوال الأساس.

\end{document}
